% Terminal apparatus of the concise study: the scope note and the References
% for the sources this companion actually uses. The revision timestamp and the
% rights colophon follow from the shared generation record.

\section{Appendix: Scope and Qualifications}

{\footnotesize
\setlength{\parskip}{0pt}
\textbf{Edition, formulary and occurrence.} This concise study treats the Mass
of the Twenty-eighth Sunday in Ordinary Time, Year~A, as the \work{Roman
Missal}, Third Edition, and the \work{Lectionary for Mass}, no.~142, give it
for the dioceses of the United States, in the national calendar alone, for the
occurrence of 11~October 2026. The antiphons and orations serve Years A, B
and~C; the readings, psalm and Alleluia verse are Year~A's. Both forms of the
Gospel and both Communion antiphons are treated: a claim that rests on
Matthew 22:11--14 does not hold for a celebration that uses the shorter form,
and a claim that rests on one Communion antiphon holds only where it is
chosen.

\textbf{Relation to the expansive study.} Every claim here is drawn from the
canonical study of this same formulary and its research records. The study
prints the appointed Scripture, sets each element in its literary and textual
setting, and develops the three interpretations at length, each with its own
literal, allegorical, moral and anagogical senses and its own difficulties;
nothing evidence-dependent is added here. The four rows on the first page
orient a reader and do not replace those distillations. The second page
condenses the study's dossier, printing the Date cells with their relation
labels, alternatives and unresolved states exactly as the generated chronology
record gives them.

\textbf{Text and English.} The Missal formulary was read for the study in the
Latin of the \work{Missale Romanum}, editio typica tertia (2002), in a
digitally typeset reproduction; no United States altar book was inspected, so
the approved English orations were not read in any witness, and the Latin
source of the Entrance Antiphon's closing words is not established. The
Lectionary's boundaries rest on the national calendar for 2026, the bishops'
conference's dated page of readings and no.~142 of the Latin \work{Ordo
lectionum Missae} (1981); the printed United States Lectionary was not
inspected, and whether it gives the psalm response as the dated page does, and
whether it admits another Alleluia verse from a common set, remain open. The
Greek of Isaiah 25 and Philippians 4:19 was not examined, and no manuscript
was. All Scripture is quoted in the Douay--Rheims version as revised by
Bishop Challoner, a public-domain study translation of the Clementine Vulgate.
It is not the text proclaimed in the celebration: the controlling English is
that of the \work{Lectionary for Mass for Use in the Dioceses of the United
States of America} and, for the antiphons and orations, that of the
\work{Roman Missal}, Third Edition, neither of which is reproduced. The
orations and antiphons are identified by Latin incipit and described; the
descriptions are not translations and may not be quoted as the texts' words.
English glosses of Latin passages, including the \work{Ordo}'s three titles,
are the editor's renderings of the Latin beside them.

\textbf{Reception and its limits.} The witnesses named are those the study's
bounded search read at the loci in the References. Chrysostom, Augustine's
sermons, his expositions of Psalms 33 and 129 and his homilies on the First
Epistle of John, the \work{Mystagogical Catecheses} and Irenaeus were read in
the English of the Nicene and Ante-Nicene Fathers, Irenaeus in a machine-read
text of the printed volume; Aquinas's \work{Summa}~I in the English Dominican
translation; Bellarmine in O'Sullivan's abridged English of 1866, the Latin
not compared; Schuster in a text layer of the 1927 English. Jerome on Matthew
and Isaiah, Hilary on Matthew, and Aquinas on Matthew, Isaiah, Ephesians and
Philippians were read on page images of the printings named; Gregory,
Aquinas's commentary on the Psalms and \work{Summa} I-II in electronic Latin
texts. No witness was collated with a modern critical edition. Nothing is
said here about who first proposed any reading, about any reading being the
Fathers' common view, or about the order in time or the mutual dependence of
the witnesses. No commentary on the Prayer over the Offerings or the Prayer
after Communion was found. The \work{Mystagogical Catecheses} are cited as
transmitted under Cyril of Jerusalem's name; their editor reports manuscripts
that give them to John of Jerusalem.

\textbf{Relations among the elements.} The liturgical books state four
relations: the first reading is chosen in relation to the Gospel
(\work{Praenotanda} 106), the psalm responds to the first reading
(\work{General Instruction of the Roman Missal} 61), the acclamation greets
the Gospel (62), and the second reading and the Gospel each belong to a
semi-continuous course. The Missal's antiphons and orations serve all three
years. Among the witnesses cited, Aquinas joins Isaiah 25:6 to Matthew 22:4
and to Psalm 22:5, and reads Psalm 22:6 as the subsequent grace whose pair of
terms the Collect uses, without citing the Collect; Augustine joins Psalm
33:11 to John 6:51. Every other connection drawn between one element and
another, and each interpretation's reading of the whole, is the editor's
synthesis, grounded in a witness's statement about his own passage where one
is cited and otherwise resting on the wording of the texts. No witness cited
comments on this Mass.

\textbf{The Council's declaration.} \work{Nostra aetate} 4 was read in the
English and Latin of the Holy See's website. It is cited as the Church's
teaching on how the identification of the first invited with a people is
presented and applied, not as a judgment on the exegesis of any Father.

\textbf{Chronology.} The Date cells print the shared chronology corpus's
answers with their relations and alternatives unchanged. The first reading's
answer dates the prophet's ministry in one author's account, not the oracle.
The Psalter's common composition bound is not a date for any of the three
psalms. The superscription answer of Psalm 33 dates the episode its title
names and rests on Haydock's notes, whose pages were not read. The Gospel's
event answer dates the telling of the parable in one author's chronology; its
composition answers, and those of the epistles, are dates of the books.

\textbf{Rights.} The Douay--Rheims, the Clementine Vulgate, the Nicene and
Post-Nicene and Ante-Nicene Fathers, O'Sullivan's Bellarmine, the English
Dominican \work{Summa}, the volumes of Migne, the Venice, Parma and Li\`ege
printings of Aquinas, Schuster's English and the \work{Catholic Encyclopedia}
are in the public domain. Gregory's Latin is quoted in short passages from a
licensed electronic edition, and the Corpus Thomisticum texts of Aquinas in
short phrases. The Latin and English of the Missal, the \work{Ordo lectionum
Missae}, the Lectionary, the \work{General Instruction}, the \work{Nova
Vulgata}, \work{Nostra aetate} and the introduction to the Psalms of the New
American Bible are protected; they are cited, and only incipits, the three
titles of the \work{Ordo}, short phrases of the Council's declaration and a
few words needed by the argument are quoted.
\par}

\phantomsection
\section*{References}
% A starred heading sets no mark, so without this the Appendix above would go
% on heading the pages of the bibliography. The heading stays starred because
% check-content-preflight finds the References section by that exact form.
\runninghead{References}
\addcontentsline{toc}{section}{References}

{\footnotesize
\subsection*{Liturgical books, Bibles and texts}
\begin{itemize}[itemsep=0pt,topsep=.2em]
\item \work{Missale Romanum}, editio typica tertia (Typis Vaticanis, 2002), \latin{Dominica XXVIII ``per annum''}.
\item \work{Ordo lectionum Missae}, editio typica altera (Libreria Editrice Vaticana, 1981): \work{Praenotanda} 80, 106; no.~142, Year~A.
\item \work{Lectionary for Mass for Use in the Dioceses of the United States of America}, second typical edition, vol.~I, no.~142, as witnessed by the United States Conference of Catholic Bishops, \work{Liturgical Calendar for the Dioceses of the United States of America} 2026, and its page of readings for 11~October 2026, \url{https://bible.usccb.org/bible/readings/101126.cfm}.
\item \work{General Instruction of the Roman Missal}, United States edition of 2011 with the emendations of 2021, nos.~61--62.
\item The Holy Bible, Douay--Rheims version, revised by Bishop Richard Challoner (Project Gutenberg text): 1~Kgs 21:10--15; Ps 16:15; 22; 33:1, 11; 129:3--4; Is 1:1; 24:23--25:10; Mt 7:13--14; 21:23--22:14; Rom 11:29; 1~Cor 10:5; 11:26; 15:54; 2~Cor 5:4; Eph 1:1, 17--18; 3:1; Phil 1:1, 7, 13; 4:10--22; 1~Tim 1:5; 1~Pet 3:21; 2~Pet 1:4; 1~Jn 3:2--3.
\item \work{Biblia Sacra Vulgatae editionis} (Clementine Vulgate), eBible.org text: Is 25:6--8; Mt 22:9; Phil 4:13, 19; 1~Pet 3:21.
\item \work{Nova Vulgata Bibliorum Sacrorum editio}, as published by the Holy See at vatican.va: Phil 4:19.
\end{itemize}

\subsection*{Fathers, Doctors and saints}
\begin{itemize}[itemsep=0pt,topsep=.2em]
\item St Irenaeus of Lyons, \work{Against Heresies} IV.36.5--6, Ante-Nicene Fathers, vol.~1 (Buffalo, 1887), in a machine-read text.
\item St Hilary of Poitiers, \work{Commentarius in Matthaeum} 22.3--7 (Migne, PL 9, Paris, 1844, cols.~1042--1044).
\item St John Chrysostom, \work{Homilies on Matthew} 69.1--2, Nicene and Post-Nicene Fathers, 1st ser., vol.~10; \work{Homilies on Ephesians} 3 and \work{Homilies on Philippians} 15, 1st ser., vol.~13.
\item The \work{Mystagogical Catecheses} ascribed to St Cyril of Jerusalem, 4.3, 4.7, Nicene and Post-Nicene Fathers, 2nd ser., vol.~7.
\item St Jerome, \work{Commentarii in Matthaeum} III, on 22:1--14 (Migne, PL 26, Paris, 1845, cols.~159--161); \work{Commentarii in Isaiam} VIII, on 25:6--8 (Migne, PL 24, Paris, 1865, col.~300).
\item St Augustine, \work{Sermo} 90.1, 4--6 and 95.5, 7, Nicene and Post-Nicene Fathers, 1st ser., vol.~6; \work{Enarrationes in Psalmos} 33 and 129.1--3, 1st ser., vol.~8; \work{In epistolam Iohannis ad Parthos tractatus} 4.5--7, 1st ser., vol.~7.
\item St Gregory the Great, \work{Homiliae in Evangelia} 38.1, 3, 5--12, 14, 16, Latin (Wikisource text).
\item St Thomas Aquinas, \work{Super Evangelium Matthaei} c.~22, in \work{Opera}, editio altera Veneta, vol.~3 (Venice: Bettinelli, 1745), pp.~279--282; \work{Expositio super Isaiam ad litteram} c.~25, in \work{Opera omnia}, vol.~14 (Parma, 1863), pp.~501--502; \work{Super epistolam ad Ephesios} c.~1, lect.~6, and \work{Super epistolam ad Philippenses} c.~4, lect.~2, in \work{In omnes S.~Pauli Apostoli epistolas commentaria}, vol.~2 (Li\`ege: Dessain, 1857), pp.~272 and 401--403; \work{Postilla super Psalmos} 22 (Corpus Thomisticum text); \work{Summa theologiae} I, q.~12, a.~6 (English Dominican translation), and I-II, q.~110, a.~3, and q.~111, a.~3 (Corpus Thomisticum text).
\item St Robert Bellarmine, \work{A Commentary on the Book of Psalms}, trans. J. O'Sullivan (1866, abridged), in the eCatholic2000 transcription: on Ps~33:10--11 and Ps~129:3--5.
\end{itemize}

\subsection*{Liturgical commentary and Magisterium}
\begin{itemize}[itemsep=0pt,topsep=.2em]
\item Blessed Ildefonso Schuster, \work{The Sacramentary (Liber Sacramentorum)}, vol.~3 (London, 1927), p.~173.
\item Second Vatican Council, Declaration \work{Nostra aetate} (28~October 1965), no.~4, English and Latin texts at vatican.va.
\end{itemize}

\subsection*{Chronology sources}
\begin{itemize}[itemsep=0pt,topsep=.2em]
\item \work{The Catholic Encyclopedia} (New York, 1908--1912): Corbett, ``David'' (vol.~4, 1908); ``Epistle to the Ephesians'' (vol.~5, 1909); Souvay, ``Isaias'', Maas, ``Jesus Christ'', and Drum, ``Epistles of Saint John'' (vol.~8, 1910); Jacquier, ``Gospel of St.~Matthew'' (vol.~10, 1911); Prat, ``St.~Paul'' (vol.~11, 1911); Vander Heeren, ``Epistle to the Philippians'' (vol.~12, 1911); Durand, ``The New Testament'' (vol.~14, 1912); read in the New Advent text.
\item United States Conference of Catholic Bishops, \work{New American Bible, Revised Edition}, introduction to the Psalms, \url{https://bible.usccb.org/bible/psalms/0} (summarized; protected).
\item George Leo Haydock, notes on 1~Kings 21 and Psalm~70, in the Douay--Rheims Bible with Haydock's commentary, as the chronology corpus records them for the A.M.\ figure.
\end{itemize}
\par}
